\section{Đa cộng tuyến và Hồi quy Ridge}

\subsection{Khái niệm Đa cộng tuyến (Multicollinearity)}

\textbf{Định nghĩa:} Đa cộng tuyến là hiện tượng xảy ra trong mô hình hồi quy bội khi hai hay nhiều biến độc lập (đặc trưng) có sự tương quan tuyến tính chặt chẽ với nhau. Điều này có nghĩa là một biến độc lập có thể được dự đoán với độ chính xác cao thông qua các biến độc lập khác trong cùng một mô hình.

\textbf{Hậu quả đối với mô hình OLS:}
\begin{itemize}
    \item \textbf{Ma trận gần suy biến:} Khi xảy ra đa cộng tuyến hoàn hảo hoặc gần hoàn hảo, các cột của ma trận thiết kế $X$ không còn độc lập tuyến tính. Hậu quả là định thức của ma trận $(X^T X)$ tiến gần về 0, làm cho phép nghịch đảo $(X^T X)^{-1}$ trong công thức nghiệm OLS trở nên cực kỳ thiếu ổn định hoặc không thể thực hiện được.
    \item \textbf{Trọng số bị thổi phồng và sai số lớn:} Các hệ số hồi quy ước lượng ($\hat{\beta}$) sẽ có giá trị tuyệt đối rất lớn và phương sai cực cao. Mô hình trở nên nhạy cảm quá mức: chỉ một thay đổi nhỏ trong tập dữ liệu huấn luyện cũng có thể làm thay đổi hoàn toàn dấu và độ lớn của các hệ số $\beta$.
    \item \textbf{Mất khả năng diễn giải:} Rất khó để đánh giá chính xác tác động độc lập của từng biến đối với biến mục tiêu $y$, làm sai lệch kết quả của các kiểm định giả thuyết (như t-test).
\end{itemize}

\textbf{Cách phát hiện bằng Hệ số phóng đại phương sai (VIF):}
Để phát hiện đa cộng tuyến, phương pháp phổ biến nhất là sử dụng Hệ số phóng đại phương sai (Variance Inflation Factor - VIF) cho từng biến $j$. Công thức được định nghĩa như sau:

\begin{equation}
    VIF_j = \frac{1}{1 - R_j^2}
\end{equation}

Trong đó, $R_j^2$ là hệ số xác định ($R$-squared) thu được khi sử dụng chính biến $X_j$ làm biến phụ thuộc và tiến hành hồi quy nó theo tất cả các biến độc lập còn lại. 
\begin{itemize}
    \item Nếu $VIF_j = 1$: Biến $X_j$ không có tương quan với các biến khác.
    \item Nếu $VIF_j > 10$ (hoặc $R_j^2 > 0.9$): Đây là dấu hiệu cảnh báo hiện tượng đa cộng tuyến đang ở mức độ nghiêm trọng và cần được xử lý.
\end{itemize}

\subsection{Hồi quy Ridge}

Để giải quyết vấn đề ma trận $(X^T X)$ suy biến do đa cộng tuyến, phương pháp Hồi quy Ridge được áp dụng. Ý tưởng cốt lõi của phương pháp này là chấp nhận hy sinh một chút tính "không chệch" (unbiasedness) của ước lượng OLS để đổi lấy một phương sai nhỏ hơn rất nhiều, từ đó giúp mô hình tổng quát hóa tốt hơn.

\textbf{Công thức toán học của Hồi quy Ridge:} \\
Thay vì chỉ tối thiểu hóa tổng bình phương phần dư (RSS) như OLS thông thường, Hồi quy Ridge thêm vào hàm mất mát một "khoản phạt" (penalty term) tỷ lệ thuận với tổng bình phương các hệ số hồi quy (không bao gồm hệ số chặn $\beta_0$). Hàm mục tiêu mới trở thành:

\begin{equation}
    \hat{\beta}_{ridge} = \arg \min_{\beta} \left\{ ||y - X\beta||_2^2 + \lambda ||\beta||_2^2 \right\}
\end{equation}

Trong đó, $\lambda \ge 0$ là siêu tham số điều chỉnh kiểm soát mức độ phạt:
\begin{itemize}
    \item Khi $\lambda = 0$: Khoản phạt bị triệt tiêu, Hồi quy Ridge hoàn toàn trở về mô hình OLS tiêu chuẩn.
    \item Khi $\lambda \rightarrow \infty$: Khoản phạt trở nên vô cùng lớn, ép tất cả các hệ số $\beta$ (trừ $\beta_0$) tiến dần về 0.
\end{itemize}

\textbf{Công thức nghiệm đóng:} \\
Bằng cách lấy đạo hàm hàm mất mát trên theo $\beta$ và cho bằng 0, ta thu được nghiệm phân tích của Hồi quy Ridge:

\begin{equation}
    \hat{\beta}_{ridge} = (X^T X + \lambda I)^{-1} X^T y
\end{equation}

Trong đó, $I$ là ma trận đơn vị bậc $p+1$. Sự xuất hiện của thành phần $\lambda I$ được cộng trực tiếp vào đường chéo chính của $(X^T X)$ là "chìa khóa" giải quyết đa cộng tuyến. Nó đảm bảo ma trận $(X^T X + \lambda I)$ luôn luôn khả nghịch, ngay cả khi số lượng đặc trưng lớn hơn số lượng mẫu quan trắc ($p > n$) hoặc khi các biến có sự cộng tuyến hoàn hảo.

